\documentclass[12pt,a4paper]{report}

% --- PACKAGES ---
\usepackage[french]{babel}
\usepackage[T1]{fontenc}
\usepackage[utf8]{inputenc}
\usepackage{listings}
\usepackage{xcolor}
\usepackage[most]{tcolorbox}
\usepackage{geometry}

\geometry{margin=2.5cm}

% --- COULEURS ---
\definecolor{calendar-red}{rgb}{0.8, 0.1, 0.1}
\definecolor{code-bg}{rgb}{0.98, 0.98, 0.98}

% --- STYLE DES BOITES (Design corrigé) ---
\newtcolorbox{datebox}[1]{
    enhanced,
    attach boxed title to top left={yshift=-2mm, xshift=2mm},
    colback=white,
    colframe=calendar-red,
    fonttitle=\bfseries,
    colbacktitle=calendar-red,
    title=#1,
    boxed title style={size=small, sharp corners},
    drop shadow,
    bottom=2mm
}

% --- CONFIGURATION C ---
\lstset{
    language=C,
    backgroundcolor=\color{code-bg},
    basicstyle=\ttfamily\footnotesize,
    keywordstyle=\color{blue},
    commentstyle=\color{gray},
    stringstyle=\color{purple},
    breaklines=true,
    frame=none,
    numbers=left,
    numberstyle=\tiny\color{gray}
}

\begin{document}

\chapter{Module de Sélection de Date (Calendar)}

\section{Introduction}
Dans le cadre du projet de planificateur de transport intelligent au Maroc, l'utilisateur doit pouvoir sélectionner une date précise. Ce module encapsule le widget \texttt{GtkCalendar} pour offrir une interface de sélection intuitive et robuste.

\section{Définition de la Configuration}

\begin{datebox}{Structure : calendar\_config}
Cette structure permet de configurer l'affichage du calendrier (noms des jours, numéros de semaines) et de définir une date sélectionnée par défaut via le type \texttt{GDateTime}.
\end{datebox}

\begin{lstlisting}
typedef struct {
    GDateTime *selected_date;
    bool show_day_names;
    bool show_week_numbers;
    bool no_month_change;
    widget_style_config style;
    void (*on_day_selected)(GtkCalendar *calendar, gpointer user_data);
    gpointer user_data;
} calendar_config;
\end{lstlisting}

\section{Macro d'Initialisation}

\begin{datebox}{Macro : CALENDAR\_CONFIG\_DEFAULT}
Comme pour les autres modules, cette macro garantit que le calendrier s'affiche avec des paramètres standard (noms des jours activés, changement de mois autorisé) dès sa création.
\end{datebox}

\begin{lstlisting}
#define CALENDAR_CONFIG_DEFAULT ((calendar_config){ \
    .selected_date = NULL, \
    .show_day_names = true, \
    .show_week_numbers = false, \
    .no_month_change = false, \
    .style = WIDGET_STYLE_CONFIG_DEFAULT, \
    .on_day_selected = NULL, \
    .user_data = NULL, \
})
\end{lstlisting}

\section{Implémentation Technique}

\begin{datebox}{Fonction : create\_calendar}
Cette fonction instancie le calendrier. Elle utilise une directive de préprocesseur \texttt{\#if GTK\_CHECK\_VERSION} pour assurer la compatibilité entre les différentes versions de GTK4 (4.20+ vs versions antérieures).
\end{datebox}

\begin{lstlisting}
GtkWidget *create_calendar(const calendar_config *config) {
    GtkWidget *calendar = gtk_calendar_new();
    if (config == NULL) return calendar;

    // Selection de la date avec gestion de version
    if (config->selected_date != NULL) {
#if GTK_CHECK_VERSION(4, 20, 0)
        gtk_calendar_set_date(GTK_CALENDAR(calendar), config->selected_date);
#else
        gtk_calendar_select_day(GTK_CALENDAR(calendar), config->selected_date);
#endif
    }
    
    // Connexion du signal de selection
    if (config->on_day_selected != NULL) {
        g_signal_connect(calendar, "day-selected", G_CALLBACK(config->on_day_selected), config->user_data);
    }

    apply_widget_style(calendar, &config->style);
    return calendar;
}
\end{lstlisting}

\end{document}